\scanpage{14}
- \mspage{11} -
6.\quad \uline{Applications}.

\uline{Proposition} 6.1.\quad \uline{Si} $E$ \uline{est une courbe de genre un sur un corps algébriquement clos}, $E$ \uline{est de type elliptique} (\uline{resp. multiplicatif}, \uline{resp. additif}) \uline{si et seulement si} $\Delta\ne0$ (\uline{resp.} $\Delta=0$ et $c_4\ne0$, \uline{resp.} $\Delta=c_4=0$).

a) Lorsque la caractéristique $p\ne2$, la première assertion résulte de la troisième formule (3.2). Lorsque $p=2$ et $c_4\ne0$, on vérifie aussitôt que la différentielle de (4.1) ne s'annule qu'en $(0,0)$, et $\Delta$ est elliptique si et seulement si $\Delta=a_6\ne0$. Pour $p=2$ et $c_4=0$, on raisonne de même sur (6.2).

b) Lorsque $p\ne2$, il résulte de (3.2) que $\Delta=c_4=0$ si et seulement si le deuxième membre de (3.1) a toutes ses racines égales, ce qui caractérise le cas additif.

c) Lorsque $p=2$, $\Delta=0$, $c_4\ne0$, la partie principale de (4.1) en $(0,0)$ est $y^2+xy+a_2x^2=0$, ce qui correspond à un point double à tangentes distinctes.

d) Lorsque $p=2$, $\Delta=0$, $c_4=0$, la partie principale de (6.2) en $(0,0)$ est $y^2+a_4x+a_6=0$, ce qui correspond à une singularité cuspidale.

Soit alors $E$ une courbe de genre un sur une base $S$, sans fibre de type additif. Puisque $\Delta$ et $c_4$ ne sont jamais simultanément nuls, et sont respectivement des sections de $\omega^{\otimes12}$ et $\omega^{\otimes4}$, on définit une section $j$ de $P_1/S$ par la formule
\margin{13}
